\documentclass[10pt,twocolumn]{article}

\usepackage[margin=0.72in]{geometry}
\usepackage{amsmath,amssymb,amsthm,mathtools}
\usepackage{microtype}
\usepackage{times}
\usepackage[hidelinks]{hyperref}
\usepackage{fontawesome5}
\usepackage{enumitem}
\usepackage{graphicx}

\setlist{nosep,leftmargin=*}
\setlength{\columnsep}{0.22in}
\setlength{\parindent}{1em}
\setlength{\parskip}{0pt}

\newtheorem{theorem}{Theorem}
\newtheorem{corollary}{Corollary}
\newtheorem{proposition}{Proposition}
\newtheorem{definition}{Definition}
\newtheorem{problem}{Problem}
\newtheorem{remark}{Remark}

\newcommand{\NN}{\mathbb{N}}
\newcommand{\ZZ}{\mathbb{Z}}
\newcommand{\RR}{\mathbb{R}}
\newcommand{\CC}{\mathbb{C}}
\newcommand{\eps}{\varepsilon}
\newcommand{\argmin}{\operatorname*{arg\,min}}
\newcommand{\supp}{\operatorname{supp}}
\newcommand{\rank}{\operatorname{rank}}
\newcommand{\figurebox}[1]{%
\fbox{\begin{minipage}{0.92\columnwidth}\small #1\end{minipage}}}

\title{\vspace{-1.0em}Inverse Takeoff Kernels:\\
System Identification for Recursive Redundancy}
\author{Jonathan R. Landers}
\date{June 2026\\[0.25em]
\href{https://github.com/jonland82/finite-response-kernels}{\faGithub\ Project Repository}}

\begin{document}
\maketitle
\vspace{-1.4em}

\begin{abstract}
Velocity takeoff kernels describe how redundant work becomes visible in a
recursive computation before memoization collapses the recursion tree into a
directed acyclic graph of distinct subproblems. This paper develops the inverse
question. Given an observed redundancy takeoff curve, what can be inferred
about the recursive mechanisms that produced it? The answer is deliberately
two-stage. First, under a finite-exponential response model, the observed
takeoff profile admits a modal expansion whose poles and amplitudes can be
recovered exactly from noiseless samples by Hankel, Prony, matrix-pencil, or
minimal-recurrence methods, and approximately from noisy traces by regularized
system identification. A dominance theorem then explains when recovered modes
become visible in the composite response: the largest-radius mode controls the
tail, while pairwise crossover estimates locate the handoff between short-lived
and slow-settling components. Second, those response modes constrain but do not
generally identify literal branch sources. A combined recurrence may be split
into fast and slow channels in many observationally equivalent ways. The
resulting principle is that takeoff identifies response modes before it
identifies source mechanisms. Source-level reconstruction becomes meaningful
only after adding structure such as bounded lag, nonnegative integer branch
counts, sparsity, separated supports, or multiple instrumented traces.
\end{abstract}

\section{Introduction}

The forward theory of takeoff kernels starts with a recursive schema and asks
what redundancy response it generates. In a finite-lag model,
\[
        N_n = 1+\sum_{j\in J} a_j N_{n-j},
\]
the tree velocity
\[
        u_n=\log N_{n+1}-\log N_n
\]
approaches a terminal speed \(\alpha\). The normalized profile
\[
        F_n=\frac{u_n}{\alpha}
\]
is a step response, and the takeoff kernel
\[
        \kappa_n=F_n-F_{n-1}, \qquad F_{-1}=0,
\]
records the finite shape of redundancy onset. The subdominant roots of the
recurrence determine whether the response is immediate, delayed, monotone,
oscillatory, mixed, or slowly settling.

That forward picture suggests a sharper empirical question. If a recursive
program is instrumented and its redundant work is observed, can the response be
read backward? A trace of \(N_n\) is not just a cost curve. After normalization,
it is a finite response shape. It may contain a fast initial component, a slow
tail, a small alternating ripple, or a near-periodic settling mode. The inverse
problem asks how much of that structure can be recovered from the curve itself:
\[
        \begin{gathered}
        \text{observed takeoff}
        \longrightarrow
        \text{modal decomposition}
        \\
        \longrightarrow
        \text{candidate recurrence mechanisms}.
        \end{gathered}
\]
The analogy is system identification, but the object being identified is not a
physical plant. It is the redundant response of a recursion tree before the
memoized quotient removes repeated labels. The observed profile is therefore
closer to a measured impulse or step response than to a conventional runtime
bound.

The main caution is identifiability. A composite recurrence may have two
channels,
\[
        a_j=b_j+c_j,
\]
where \(b_j\) is a fast channel and \(c_j\) is a slow channel. The observed tree
sequence depends on the combined coefficients \(a_j\), not on the chosen split.
Thus a takeoff trace can identify response modes more directly than it can
identify literal source branches. This is the paper's organizing distinction:
modal structure is the primary observable; source structure is a constrained
interpretation layered on top of it.

\paragraph{Proposed contributions.}
The paper develops six pieces.
\begin{enumerate}
    \item Define inverse takeoff as a recovery problem for observed redundancy
    velocity profiles.
    \item Prove an exact modal recovery theorem in the noiseless
    finite-exponential case.
    \item Characterize when recovered modes become apparent in the composite
    response through tail dominance and crossover estimates.
    \item Prove a nonidentifiability result for latent branch-channel
    decompositions.
    \item Develop constrained reconstruction criteria for plausible recursive
    mechanisms.
    \item Demonstrate the method on synthetic and real recursive traces.
\end{enumerate}

\begin{figure}[t]
\centering
\figurebox{
\textbf{Planned Figure 1: inverse takeoff pipeline.}
Show the measured curve \(N_n\), the normalized residual
\(Y_n=F_n-1\), the Hankel singular spectrum used to choose modal order,
the recovered modes \(\theta_r\), and the final constrained recurrence
candidates. The point of the figure is to separate response identification
from source interpretation.
}
\caption{Conceptual pipeline for inverse takeoff.}
\label{fig:pipeline-placeholder}
\end{figure}

\section{Observed Takeoff Data}

Let \(N_n\) be the number of calls in the naive recursion tree and \(M_n\) the
number of distinct subproblem labels under memoization. Define
\[
        u_n=\log N_{n+1}-\log N_n,
        \qquad
        w_n=\log M_{n+1}-\log M_n,
\]
and
\[
        v_n=u_n-w_n.
\]
The sequence \(u_n\) is the tree-side velocity; it records how quickly
recursive multiplicity is being created. The sequence \(w_n\) is the
state-space velocity; it records how quickly the memoized DAG itself is
growing. The overlap velocity \(v_n\) is their difference. This distinction is
important in the inverse setting because two different response curves can
produce a visually simple difference.

When \(u_n\to\alpha>0\), the tree takeoff profile is
\[
        F_n=\frac{u_n}{\alpha}.
\]
When \(w_n\to\beta>0\), the DAG profile is
\[
        G_n=\frac{w_n}{\beta}.
\]
The overlap velocity obeys
\[
        v_n-(\alpha-\beta)
        =
        \alpha(F_n-1)-\beta(G_n-1).
\]
Thus, if only \(v_n\) is observed, tree and DAG modes may cancel or mask one
another. The inverse problem is best posed using \(N_n\) and \(M_n\) separately
when those traces are available.

The cleanest data product is therefore
\[
        \mathcal{D}_T
        =
        \{(N_n,M_n):0\le n\le T\}.
\]
From \(\mathcal{D}_T\), one computes \(u_n,w_n,v_n\), estimates terminal
speeds \(\alpha,\beta\), and then forms normalized residuals. If only the tree
trace is available, the main residual is
\[
        Y_n=F_n-1.
\]
If only overlap velocity is available, the measured residual is instead
\[
        Z_n
        =
        v_n-(\widehat{\alpha}-\widehat{\beta}),
\]
which should not be interpreted as a pure tree mode unless \(w_n\) is known to
settle faster or be negligible.

In measured data, write
\[
        \widehat Y_n = Y_n+\eps_n,
        \qquad n=n_0,\ldots,n_1,
\]
where \(\eps_n\) includes measurement noise, terminal-speed estimation error,
boundary effects, and model misspecification. The finite prefix is often
contaminated by base-case conventions, so the inverse procedure should allow a
warmup cutoff \(n_0\). The tail, meanwhile, may be too short to expose a slow
mode. A credible inverse analysis must therefore report the chosen observation
window, not only the fitted parameters.

\paragraph{Estimating terminal speed.}
For a tree trace, a simple estimator is a tail average
\[
        \widehat{\alpha}
        =
        \frac{1}{|\mathcal{T}|}\sum_{n\in\mathcal{T}}u_n
\]
over a stable tail window \(\mathcal{T}\). A more robust version fits
\[
        \log N_n
        =
        c+\alpha n+\delta_n
\]
on the tail and then studies the residual velocity. The inverse method should
be run under several plausible tail windows; a mode that disappears under small
window shifts is better treated as a numerical artifact than as a structural
claim.

\begin{figure}[t]
\centering
\figurebox{
\textbf{Planned Figure 2: data normalization.}
Panel (a): raw \(N_n\) on a log scale. Panel (b): tree velocity \(u_n\)
with the fitted terminal speed \(\widehat{\alpha}\). Panel (c): residual
\(Y_n=\widehat F_n-1\) after a warmup cutoff. Panel (d): sensitivity of the
residual to the chosen tail window.
}
\caption{From call counts to an inverse-takeoff residual.}
\label{fig:normalization-placeholder}
\end{figure}

\section{Three Inverse Problems}

\begin{problem}[Modal inverse takeoff]
Given samples \(Y_0,\ldots,Y_T\), recover a small modal model
\[
        F_n-1
        \approx
        \sum_{r=1}^{m} c_r\theta_r^n.
\]
For complex conjugate modes this can be written as a sum of damped sinusoids,
\[
        A_r r_r^n\cos(\omega_r n+\phi_r).
\]
The target is not merely interpolation. The fitted modes should predict
unobserved tail values, explain the observed settling envelope, and remain
stable under modest changes of the observation window.
\end{problem}

\begin{problem}[Combined recurrence reconstruction]
Given \(N_0,\ldots,N_T\), recover a bounded-lag recurrence
\[
        N_n
        \approx
        1+\sum_{j=1}^{L}a_jN_{n-j},
        \qquad
        a_j\ge 0,
\]
preferably with \(a_j\in\ZZ_{\ge0}\) when the model represents branch counts.
This problem asks for a combined mechanism. It does not yet ask which branch or
subroutine supplied each lag.
\end{problem}

\begin{problem}[Latent channel decomposition]
Given recovered combined coefficients \(a_j\), recover a decomposition
\[
        a_j=b_j+c_j
\]
into fast, slow, or otherwise interpretable channels.
\end{problem}

These problems are ordered by identifiability. Modal inverse takeoff is the
best-posed. Combined recurrence reconstruction is plausible under bounded-lag
and positivity assumptions. Latent channel decomposition is not identifiable in
general and requires additional structure.

The ordering also suggests how results should be reported. The first output is
a modal table
\[
        (\theta_r,c_r)
        =
        (r_re^{i\omega_r},c_r),
        \qquad r=1,\ldots,m,
\]
with damping \(r_r\), frequency \(\omega_r\), phase, and uncertainty. The
second output is a small set of candidate recurrences \(a\), each with a
forward-simulated takeoff curve. The final output, if attempted, is a source
story such as ``one local branch and one delayed branch.'' Only the first of
these is directly identified by a single trace.

\section{Exact Modal Recovery}

Suppose the takeoff tail has the finite-exponential form
\[
        y_n:=F_n-1=\sum_{r=1}^{m}c_r\theta_r^n,
\]
with distinct nonzero modes \(\theta_r\) and nonzero amplitudes \(c_r\).
Define the Hankel matrix
\[
        H^{(s)}_{ij}=y_{s+i+j},
        \qquad 0\le i,j\le q-1.
\]
Let \(V_q\) be the \(q\times m\) Vandermonde matrix
\[
        (V_q)_{ir}=\theta_r^i,
        \qquad 0\le i\le q-1,\quad 1\le r\le m.
\]
Then
\[
        H^{(s)}
        =
        V_q
        \operatorname{diag}(c_1\theta_1^s,\ldots,c_m\theta_m^s)
        V_q^{T}.
\]
Consequently, for \(q\ge m\), the infinite noiseless Hankel matrix has rank
\(m\). This is the algebraic signature of a finite modal response: the data
live on an \(m\)-dimensional recurrence model even though the observed trace may
look like a composite curve.

\begin{theorem}[Noiseless modal recovery]
Assume
\[
        y_n=\sum_{r=1}^{m}c_r\theta_r^n
\]
with distinct nonzero \(\theta_r\) and nonzero \(c_r\). Then the \(2m\)
consecutive samples \(y_0,\ldots,y_{2m-1}\) determine the monic annihilating
polynomial
\[
        p(t)=t^m+p_{m-1}t^{m-1}+\cdots+p_0
        =
        \prod_{r=1}^{m}(t-\theta_r),
\]
uniquely. The modes \(\theta_r\) are the roots of \(p\), and the amplitudes
\(c_r\) are recovered uniquely by a Vandermonde system.
\end{theorem}

\begin{proof}
The sequence satisfies
\[
        y_{n+m}+p_{m-1}y_{n+m-1}+\cdots+p_0y_n=0
\]
for every \(n\), because each \(\theta_r\) is a root of \(p\). For
\(n=0,\ldots,m-1\), this gives the Hankel system
\[
        \begin{bmatrix}
        y_0 & y_1 & \cdots & y_{m-1}\\
        y_1 & y_2 & \cdots & y_m\\
        \vdots & \vdots & \ddots & \vdots\\
        y_{m-1} & y_m & \cdots & y_{2m-2}
        \end{bmatrix}
        \begin{bmatrix}
        p_0\\ p_1\\ \vdots\\ p_{m-1}
        \end{bmatrix}
        =
        -
        \begin{bmatrix}
        y_m\\ y_{m+1}\\ \vdots\\ y_{2m-1}
        \end{bmatrix}.
\]
The coefficient matrix is \(V_m\operatorname{diag}(c_r)V_m^T\), hence is
nonsingular because the \(\theta_r\) are distinct and the \(c_r\) are nonzero.
Thus the polynomial coefficients are determined. Once \(p\) is known, its roots
give the modes. The first \(m\) equations
\[
        y_n=\sum_{r=1}^{m}c_r\theta_r^n,
        \qquad 0\le n<m,
\]
form a Vandermonde system for the amplitudes.
\end{proof}

This theorem is classical Prony theory in the notation of takeoff kernels. In
practice one uses a numerically stable variant: matrix pencil, ESPRIT,
Hankel-SVD, or prediction-error fitting.

\begin{corollary}[Minimal recurrence]
Under the hypotheses of the theorem, \(y_n\) satisfies a unique minimal linear
recurrence of order \(m\),
\[
        y_{n+m}
        =
        -\sum_{\ell=0}^{m-1}p_\ell y_{n+\ell},
\]
and no recurrence of smaller order generates the same infinite sequence.
\end{corollary}

\begin{proof}
The recurrence follows from the annihilating polynomial. If a smaller
recurrence existed, its characteristic polynomial would annihilate every
\(\theta_r\), but a nonzero polynomial of degree less than \(m\) cannot have
the \(m\) distinct roots \(\theta_1,\ldots,\theta_m\).
\end{proof}

\begin{proposition}[Tail dominance and modal visibility]
Let
\[
        y_n=\sum_{r=1}^{m}c_r\theta_r^n,
\]
where the modes \(\theta_r\) are distinct and the amplitudes \(c_r\) are
nonzero. Suppose the modes are ordered so that
\[
        |\theta_1|>\lambda,
        \qquad
        \lambda=\max_{r\ge2}|\theta_r|.
\]
Then the largest-radius mode is asymptotically visible:
\[
        y_n
        =
        c_1\theta_1^n
        \left(
        1+O\left(\left(\frac{\lambda}{|\theta_1|}\right)^n\right)
        \right).
\]
More generally, for two modes \(r\) and \(s\) with
\(|\theta_r|>|\theta_s|\), mode \(r\) dominates mode \(s\) by a factor
\(\eta>1\) once
\[
        n
        \ge
        \frac{\log(\eta |c_s|/|c_r|)}
             {\log(|\theta_r|/|\theta_s|)}.
\]
When equality is used with \(\eta=1\), this gives the formal crossover index
\[
        n_{r,s}
        =
        \frac{\log(|c_s|/|c_r|)}
             {\log(|\theta_r|/|\theta_s|)}.
\]
\end{proposition}

\begin{proof}
Factor out the largest-radius term:
\[
        y_n
        =
        c_1\theta_1^n
        \left(
        1+\sum_{r=2}^{m}
        \frac{c_r}{c_1}\left(\frac{\theta_r}{\theta_1}\right)^n
        \right).
\]
The parenthesized correction is
\[
        1+O\left(\left(\frac{\lambda}{|\theta_1|}\right)^n\right),
\]
because \(|\theta_r/\theta_1|\le \lambda/|\theta_1|<1\). For the pairwise
visibility estimate, solve
\[
        |c_r|\,|\theta_r|^n
        \ge
        \eta |c_s|\,|\theta_s|^n
\]
for \(n\). This gives the displayed threshold.
\end{proof}

\begin{remark}
The proposition is an asymptotic and pairwise visibility statement. In a finite
window, a mode can still be hidden by boundary effects, noise, nearby modes,
or cancellation among several comparable components. This is why the numerical
procedure should report Hankel singular values, window sensitivity, and
holdout prediction error rather than only a list of fitted modes.
\end{remark}

\paragraph{Kernel visibility.}
The same mode contributes to the takeoff kernel by finite difference:
\[
        \kappa_n
        =
        F_n-F_{n-1}
        \quad\leadsto\quad
        c_r(\theta_r-1)\theta_r^{n-1}.
\]
Thus very slow modes with \(\theta_r\approx1\) may be subtle in the kernel
while still controlling the long settling of the profile \(F_n\).

\paragraph{Oscillatory modes.}
If a conjugate pair has
\[
        \theta,\overline{\theta}=re^{\pm i\omega},
\]
then its real contribution can be written
\[
        C r^n\cos(\omega n+\phi).
\]
Thus the fitted parameters have direct qualitative meaning: \(r\) is the
settling rate, \(\omega\) is the ringing frequency, and \(\phi\) is the phase
set by the boundary convention and early response. This is why the same modal
language covers monotone smoothing, alternating overshoot, and near-periodic
takeoff.

\begin{figure}[t]
\centering
\figurebox{
\textbf{Planned Figure 3: modal visibility.}
Plot two or three components \(|c_r||\theta_r|^n\) and their sum. Mark the
crossover index \(n_{r,s}\) where a weaker but slower-decaying component
overtakes a stronger startup component. A second panel should show the same
effect in the differenced kernel \(\kappa_n\), where modes near
\(\theta=1\) are attenuated by \(\theta-1\).
}
\caption{When response modes become visible in a composite takeoff curve.}
\label{fig:visibility-placeholder}
\end{figure}

\paragraph{From modes to recurrence roots.}
For a finite-lag recurrence with denominator
\[
        Q(z)=1-\sum_{j=1}^{L}a_jz^j,
\]
let \(\rho=e^{-\alpha}\) be the dominant positive root. A non-dominant
singularity \(\zeta_r\) contributes a mode
\[
        \theta_r=\frac{\rho}{\zeta_r}.
\]
Thus modal recovery gives candidate singularities
\[
        \widehat{\zeta}_r=\frac{\widehat{\rho}}{\widehat{\theta}_r}.
\]
If enough roots are recovered and the bounded-lag model is correct, these
roots constrain \(Q\), and therefore the combined coefficients \(a_j\).

One reconstruction route is root-based. Fix a maximum lag \(L\) and write
\[
        Q_L(z)=1-\sum_{j=1}^{L}a_jz^j.
\]
The ideal constraints are
\[
        Q_L(\rho)=0,
        \qquad
        Q_L(\zeta_r)=0
        \quad \text{for recovered singularities } \zeta_r.
\]
With noisy modes, replace them by a weighted residual:
\[
        \min_{a\in\RR^L}
        \left|1-\sum_{j=1}^{L}a_j\widehat\rho^{\,j}\right|^2
        +
        \sum_r \omega_r
        \left|1-\sum_{j=1}^{L}a_j\widehat\zeta_r^{\,j}\right|^2.
\]
The weights \(\omega_r\) should be smaller for unstable or weakly recovered
modes. The solution is only a candidate denominator; it must still be checked
against the original sequence and against structural constraints such as
\(a_j\ge0\) or \(a_j\in\ZZ_{\ge0}\).

A second route avoids root reconstruction and fits the recurrence directly:
\[
        a^{(L)}
        =
        \argmin_{a_1,\ldots,a_L}
        \sum_{n=L}^{T}
        \left(
            N_n-1-\sum_{j=1}^{L}a_jN_{n-j}
        \right)^2.
\]
This direct fit is usually more stable numerically. The modal model then acts
as an explanation and diagnostic for the fitted recurrence: its roots should
produce the same observed damping, ringing, and crossover behavior.

\paragraph{Repeated roots and polynomial factors.}
If \(Q\) has repeated roots, the modal expansion contains terms
\[
        n^\ell \theta^n.
\]
The first version of the follow-up should state the simple-root theory first,
then add repeated roots as a standard extension.

\section{Nonidentifiability of Sources}

\begin{proposition}[Branch-channel decomposition is not identifiable]
The latent split \(a_j=b_j+c_j\) cannot be recovered from the tree sequence
\((N_n)\) without additional assumptions. Any two splits with the same
combined coefficients \(a_j\) generate the same recurrence, the same tree
sequence, and the same takeoff kernel.
\end{proposition}

\begin{proof}
The recurrence depends on \(b_j\) and \(c_j\) only through their sum:
\[
        N_n
        =
        1+\sum_j (b_j+c_j)N_{n-j}
        =
        1+\sum_j a_jN_{n-j}.
\]
Therefore all splits with the same \(a_j\) are observationally equivalent from
the perspective of \(N_n\), \(u_n\), \(F_n\), and \(\kappa_n\).
\end{proof}

Equivalently, a single trace identifies at most an equivalence class
\[
        [(b,c)]
        =
        \{(b',c'):\ b'_j+c'_j=a_j\ \text{for all }j\}.
\]
Every element of this class has the same response. The labels ``fast'' and
``slow'' become meaningful only when they are tied to information not present
in the single aggregate trace: separated lag supports, known code paths,
interventions that disable one channel, or multiple measurements in which the
channels are perturbed differently.

A simple example is
\[
        a_1=2,\qquad a_4=2.
\]
The split
\[
        b_1=2,\ b_4=0,
        \qquad
        c_1=0,\ c_4=2
\]
and the split
\[
        b_1=1,\ b_4=1,
        \qquad
        c_1=1,\ c_4=1
\]
produce the same combined recurrence. One interpretation is a fast channel plus
a slow channel; the other is two identical mixed channels. A single composite
takeoff trace cannot choose between them.

\begin{remark}
This is not a defect of the framework. It is the useful warning. Takeoff
identifies response modes before it identifies source mechanisms.
\end{remark}

\begin{figure}[t]
\centering
\figurebox{
\textbf{Planned Figure 4: nonidentifiability.}
Show two different channel diagrams with the same combined coefficients
\(a_1=2,a_4=2\). Below them, plot the identical \(F_n\) and \(\kappa_n\).
The figure should make clear that response equivalence is exact, not a
numerical accident.
}
\caption{Different source stories can induce the same takeoff response.}
\label{fig:nonident-placeholder}
\end{figure}

\section{Constrained Structural Recovery}

To make source reconstruction meaningful, impose structure. Natural
constraints include:
\begin{enumerate}
    \item bounded lag: \(a_j=0\) for \(j>L\);
    \item nonnegative branch counts: \(a_j\ge0\);
    \item integer branch counts: \(a_j\in\ZZ_{\ge0}\);
    \item sparsity: only a few lags are active;
    \item separated supports for fast and slow channels;
    \item multiple traces from controlled perturbations of the same program.
\end{enumerate}

For a fixed maximum lag \(L\), let \(\mathcal{A}_L\) be the feasible class. For
example,
\[
        \mathcal{A}_L^{+}
        =
        \{a\in\RR^L:\ a_j\ge0\},
        \qquad
        \mathcal{A}_L^{\ZZ}
        =
        \{a\in\ZZ_{\ge0}^{L}\}.
\]
A direct combined-recurrence fit is then
\[
        \widehat a^{(L)}
        =
        \argmin_{a\in\mathcal{A}_L}
        \sum_{n=L}^{T}
        \left(
            N_n-1-\sum_{j=1}^{L}a_jN_{n-j}
        \right)^2
        +\lambda\|a\|_1,
\]
where \(\lambda\|a\|_1\) encourages sparse active lags. A model-selection score
can then penalize lag length and support size:
\[
        \begin{aligned}
        \mathrm{Score}(L,a)
        &=
        \log\left(\frac{1}{T-L+1}
        \sum_{n=L}^{T}R_n(a)^2\right)\\
        &\quad
        +\tau |\supp(a)|+\gamma L,
        \end{aligned}
\]
where
\[
        R_n(a)=N_n-1-\sum_{j=1}^{L}a_jN_{n-j}.
\]
For branch-count reconstruction, replace the nonnegative constraint by
\[
        a_j\in\ZZ_{\ge0}
\]
and solve a small mixed-integer least-squares problem or use nonnegative least
squares followed by integer polishing.

If only the normalized profile is available, reconstruct a scale-free tree
proxy by
\[
        \log \widetilde N_{n+1}
        =
        \log \widetilde N_n+\widehat{\alpha}\widehat F_n.
\]
The multiplicative scale of \(\widetilde N_n\) is arbitrary, but the velocity
profile is preserved. This is weaker than observing \(N_n\) directly and should
be treated as a diagnostic rather than a definitive reconstruction.

\paragraph{Forward validation.}
Every candidate \(\widehat a\) should be pushed back through the forward map.
Generate \(\widehat N_n\) from
\[
        \widehat N_n
        =
        1+\sum_{j=1}^{L}\widehat a_j\widehat N_{n-j}
\]
using the same boundary convention as the data, and compare the induced
\(\widehat F_n\), \(\widehat\kappa_n\), and modal roots against the observed
ones. A recurrence that fits \(N_n\) but produces the wrong takeoff shape is a
poor explanation, even if its one-step residual is small.

\paragraph{Channel recovery with separated supports.}
Assume two channels
\[
        a=b+c,
\]
with
\[
        \begin{aligned}
        \supp(b)&\subseteq \{1,\ldots,L_f\},\\
        \supp(c)&\subseteq \{L_s,\ldots,L\},\\
        L_f&<L_s.
        \end{aligned}
\]
Then a fast/slow interpretation becomes more stable: early lags must be
assigned to \(b\), late lags to \(c\), and only the gap region remains
ambiguous. This is a plausible assumption for algorithms where one mechanism
recurs locally and another returns after a large delayed dependency.

With multiple traces, the constraint can be stronger. Suppose experiments
\(e=1,\ldots,E\) change the channel weights but preserve supports:
\[
        a^{(e)}=s_e b+t_e c.
\]
If the scalars \(s_e,t_e\) are known or separately estimable and the design
matrix \((s_e,t_e)\) has rank two, then the combined coefficients across
experiments can separate \(b\) and \(c\). This is the inverse-takeoff analogue
of adding interventions or sensors: source recovery requires more than one
aggregate view.

\section{Algorithms}

\paragraph{Algorithm 1: modal inverse takeoff.}
\begin{enumerate}
    \item Measure \(N_n\) and, if relevant, \(M_n\).
    \item Estimate \(\alpha\) from a stable tail window of \(u_n\).
    \item Form \(Y_n=\widehat F_n-1\), optionally after discarding a warmup
    prefix.
    \item Choose model order \(m\) from Hankel singular values, information
    criteria, or holdout prediction error.
    \item Estimate modes using Prony, matrix pencil, ESPRIT, or Hankel-SVD.
    \item Convert complex conjugate pairs to damping, frequency, and phase:
    \[
            \theta=re^{i\omega}
            \quad\leadsto\quad
            C r^n\cos(\omega n+\phi).
    \]
    \item Validate by forward-predicting \(F_n\), \(\kappa_n\), and the
    settling envelope.
\end{enumerate}
The modal report should include the table
\[
        \left\{
        (|\theta_r|,\arg\theta_r,c_r,n_{r,s})
        \right\},
\]
where crossover entries \(n_{r,s}\) are reported only for pairs whose
amplitudes are large enough to be visible in the observation window. This
prevents the method from over-reading tiny fitted poles.

\paragraph{Algorithm 2: constrained recurrence search.}
\begin{enumerate}
    \item Pick a maximum lag \(L_{\max}\).
    \item For each \(L\le L_{\max}\), fit \(a_1,\ldots,a_L\) under
    nonnegative or integer constraints.
    \item Simulate the recovered recurrence and compare its takeoff profile to
    the observed profile.
    \item Penalize large \(L\), dense supports, negative implied modes, and
    unstable dominant roots.
    \item Report a small set of candidate recurrences rather than a single
    overconfident source mechanism.
\end{enumerate}

\paragraph{Algorithm 3: channel interpretation.}
\begin{enumerate}
    \item Start from a recovered combined recurrence \(\widehat a\).
    \item Declare the external assumptions: separated supports, known code
    paths, intervention labels, or multiple traces.
    \item Solve for channel coefficients only inside the assumed class.
    \item Report the equivalence class left unresolved by the assumptions.
\end{enumerate}
This third step is intentionally conservative. It is where explanatory stories
enter, so it should be the most explicit about assumptions.

\section{Examples and Experiments}

The first experiments should be deliberately simple. Their purpose is not to
beat existing algorithms for recurrence solving. Their purpose is to test
whether inverse takeoff recovers the right qualitative structure from finite
traces: immediate response, delay, monotone settling, oscillation, mixed-stage
response, and failure on non-exponential tails.

Each synthetic experiment should report:
\[
        \widehat m,\qquad
        \{(\widehat\theta_r,\widehat c_r)\}_{r=1}^{\widehat m},
        \qquad
        \widehat a,
\]
along with the residuals
\[
        \sum_n |Y_n-\widehat Y_n|^2,
        \qquad
        \sum_n |F_n-\widehat F_n|^2,
        \qquad
        \sum_n |\kappa_n-\widehat\kappa_n|^2.
\]
The first residual tests modal fit; the second and third test whether the
shape of redundancy onset is reproduced.

\paragraph{Fibonacci.}
The recurrence
\[
        N_n=1+N_{n-1}+N_{n-2}
\]
has a monotone leading correction under node counting and a more visible
alternating component under leaf counting. This tests sensitivity to boundary
and counting convention.

\paragraph{Delayed binary recursion.}
The recurrence
\[
        N_n=1+2N_{n-L}
\]
has dead time in the original coordinate and immediate takeoff in the coarse
coordinate \(m=n/L\). This tests whether the inverse method recognizes delay.

\paragraph{Mixed fast/slow channels.}
Use
\[
        N_n=1+b_1N_{n-1}+c_LN_{n-L}.
\]
The intended response is a fast initial rise followed by a slow settling tail.
This is the central synthetic case for the follow-up.

\paragraph{Near-periodic slow family.}
Use the family
\[
        N_n=1+(2^L-1)N_{n-L}+2N_{n-L-1}.
\]
The dominant speed is \(\log 2\), but non-dominant roots approach the dominant
circle and create a long ringing response. This tests modal resolution when
roots are close to the unit circle in normalized coordinates.

\paragraph{Grid dynamic programs as a negative control.}
Diagonal lattice paths satisfy
\[
        P_{n,n}=\binom{2n}{n}
\]
and have algebraic rather than exponential settling. A finite-exponential
inverse model should either fail gracefully or require too many modes. This
separates finite-lag recurrence takeoff from diffusive takeoff.

\begin{figure}[t]
\centering
\figurebox{
\textbf{Planned Figure 5: synthetic benchmark grid.}
Rows: Fibonacci, delayed binary, mixed fast/slow, near-periodic slow family,
and lattice-path negative control. Columns: \(F_n\), \(\kappa_n\), Hankel
singular spectrum, recovered poles in the complex plane, and forward
simulation of the best candidate recurrence.
}
\caption{Synthetic inverse-takeoff benchmark layout.}
\label{fig:synthetic-placeholder}
\end{figure}

\begin{figure}[t]
\centering
\figurebox{
\textbf{Planned Table 1: reconstruction metrics.}
For each benchmark, report true active lags, recovered active lags, modal
order, pole error, holdout error, and whether the source decomposition is
identified, underidentified, or deliberately not attempted.
}
\caption{Experimental analysis table reserved for the completed study.}
\label{fig:metrics-placeholder}
\end{figure}

\paragraph{Real recursive traces.}
Instrument recursive programs to record:
\[
        N_n,\qquad M_n,\qquad u_n,\qquad w_n,\qquad v_n.
\]
Candidate examples include memoized search, context-free parsing recurrences,
edit-distance-style dynamic programs, and recursive enumeration tasks with
changing branching factors.

For real traces, the paper should be careful about interpretation. The right
claim is not that the recovered modes reveal the source code uniquely. The
claim is that a measured recursive workload has a reproducible response shape.
The analysis should therefore compare runs across input families, random seeds,
and boundary conventions. A mode that persists across these choices is a
stronger empirical object than a mode visible in only one fitted window.

\section{Evaluation Criteria}

The follow-up should avoid claiming that a recovered modal model is the true
program mechanism. Report several levels of evidence:
\begin{enumerate}
    \item modal residual:
    \[
            \sum_n |Y_n-\widehat Y_n|^2;
    \]
    \item Hankel-rank stability under window shifts;
    \item pole stability under bootstrap or subsampling;
    \item recurrence residual on \(N_n\);
    \item holdout prediction error for later \(n\);
    \item stability of inferred support \(\supp(a)\);
    \item sensitivity to warmup cutoff and base-case convention.
\end{enumerate}

Failure modes should be part of the story: short windows, nearly colliding
modes, noisy counts, algebraic tails, hidden DAG cancellation, and multiple
structural decompositions with the same combined coefficients.

A useful reporting convention is to classify each run at three levels:
\[
        \begin{gathered}
        \text{modal fit}
        \Rightarrow
        \text{combined recurrence}\\
        \Rightarrow
        \text{source interpretation}.
        \end{gathered}
\]
The modal fit is accepted when it predicts held-out \(Y_n\) and has stable
Hankel rank. The combined recurrence is accepted when it reproduces both
\(N_n\) and the takeoff shape. The source interpretation is accepted only when
the stated assumptions make the decomposition identifiable. This three-level
reporting keeps a successful signal fit from being mistaken for a recovered
program mechanism.

\section{Related Literature Map}

\paragraph{Analytic combinatorics and recurrence asymptotics.}
The forward theory is the standard rational-generating-function story:
singularities of the generating function determine exponential growth and
subdominant correction modes. This is the basis for translating finite-lag
recurrences into takeoff modes.

\paragraph{Prony, matrix pencil, and ESPRIT.}
These methods recover sums of exponentials and damped sinusoids from samples.
They are the closest technical literature for the exact and noisy modal inverse
takeoff problem.

\paragraph{Berlekamp-Massey, Pade approximation, and minimal recurrences.}
For exact linearly recurrent sequences, one can recover a shortest recurrence
or rational generating function from finite samples. This is the algebraic
version of the inverse recurrence problem.

\paragraph{System identification.}
Prediction-error methods, subspace identification, Ho-Kalman realization, and
Hankel-SVD methods provide the engineering language: infer a small dynamical
system from an observed impulse or step response.

\paragraph{Blind source separation.}
This literature is relevant mostly as a caution. A single composite trace does
not determine independent latent sources without additional assumptions, such
as multiple sensors, independence, non-Gaussianity, interventions, or structural
constraints.

\paragraph{Sparse inverse problems.}
Sparsity and regularization give a practical path from many possible
recurrences to a small set of interpretable candidates.

\section{Claim Map for the Complete Version}

The complete paper should separate proved algebraic claims from empirical and
algorithmic claims.

\paragraph{Exact algebraic claims.}
The finite-exponential model gives exact modal recovery from \(2m\) samples,
Hankel rank \(m\), a minimal recurrence of order \(m\), tail dominance of the
largest-radius mode, and pairwise crossover estimates. These are the stable
mathematical core of the paper.

\paragraph{Structural claims.}
Recovered modes determine candidate recurrence roots by
\[
        \widehat\zeta_r=\widehat\rho/\widehat\theta_r,
\]
and bounded-lag recurrence candidates can be fit either from these roots or
directly from \(N_n\). The structural claims should be phrased as constrained
reconstruction, not as unique source recovery.

\paragraph{Negative claims.}
Branch-channel splits \(a=b+c\) are not identifiable from one aggregate trace.
This result should remain explicit because it protects the inverse theory from
overinterpretation.

\paragraph{Empirical claims.}
The experiments should show that synthetic recurrences are recovered in the
right regime, that algebraic or diffusive tails trigger failure modes, and that
real recursive traces have reproducible response shapes. Perturbation bounds
for noisy modal recovery can be quoted from matrix-pencil or
subspace-identification literature rather than reproved.

\section{Conclusion}

The inverse takeoff problem is best framed as a measured-response problem, not
as a promise to recover hidden branches uniquely. In the finite-lag setting,
the observed redundancy response can often be decomposed into modes:
\[
        F_n-1\approx\sum_r c_r\theta_r^n.
\]
Those modes explain damping, ringing, delay, and settling. They also constrain
possible recurrence mechanisms. The mode with largest radius becomes visible in
the tail, and pairwise crossover estimates describe when slower-decaying modes
overtake stronger short-lived components.

The subtle point is that this success stops one step short of source recovery.
The same combined recurrence can support several incompatible stories about
fast and slow channels. That ambiguity is not an embarrassment; it is the line
between what the trace says and what an analyst adds. The trace says which
response modes are present. Extra assumptions say which mechanisms may have
produced them.

The remaining work is empirical as much as algebraic. The algebraic core gives
Hankel rank, exact modal recovery, minimal recurrence, visibility thresholds,
and nonidentifiability. The experimental core must show when these statements
survive finite windows, noisy counts, boundary conventions, and real recursive
programs. If that evidence is supplied, inverse takeoff kernels become a
practical diagnostic: a way to read recursive redundancy as a response signal
before trying to name its sources. The guiding principle is therefore:
\[
        \boxed{
        \begin{gathered}
        \text{Takeoff identifies response modes}\\
        \text{before it identifies source mechanisms.}
        \end{gathered}}
\]

\begin{thebibliography}{9}

\bibitem{bellman1957}
R. Bellman.
\newblock \emph{Dynamic Programming}.
\newblock Princeton University Press, 1957.

\bibitem{flajolet2009}
P. Flajolet and R. Sedgewick.
\newblock \emph{Analytic Combinatorics}.
\newblock Cambridge University Press, 2009.

\bibitem{hua1990}
Y. Hua and T. K. Sarkar.
\newblock Matrix pencil method for estimating parameters of exponentially
damped/undamped sinusoids in noise.
\newblock \emph{IEEE Transactions on Acoustics, Speech, and Signal Processing},
38(5):814--824, 1990.

\bibitem{ljung1999}
L. Ljung.
\newblock \emph{System Identification: Theory for the User}.
\newblock Prentice Hall, second edition, 1999.

\bibitem{potts2010}
D. Potts and M. Tasche.
\newblock Parameter estimation for nonincreasing exponential sums by
Prony-like methods.
\newblock \emph{Linear Algebra and its Applications}, 439(4):1024--1039, 2013.

\bibitem{prony1795}
G. de Prony.
\newblock Essai experimental et analytique.
\newblock \emph{Journal de l'Ecole Polytechnique}, 1(2):24--76, 1795.

\bibitem{vanoverschee1996}
P. Van Overschee and B. De Moor.
\newblock \emph{Subspace Identification for Linear Systems}.
\newblock Kluwer Academic Publishers, 1996.

\bibitem{hyvarinen2000}
A. Hyvarinen and E. Oja.
\newblock Independent component analysis: algorithms and applications.
\newblock \emph{Neural Networks}, 13(4--5):411--430, 2000.

\end{thebibliography}

\end{document}
